\chapter{Literature Review}
\label{ch:Literature Review}

\section{Optimization Problems}
\label{sec:optimization-problems}

\subsection{Combinatorial Optimization}
\label{subsec:combinatorial-optimization}

This paper follows the conventional notation for combinatorial optimization introduced by
\citet{papadimitriou1998combinatorial}. The space of candidate solutions for a CO problem
instance $s$ is denoted by $\mathcal{X}_s = \{0, 1\}^N$, and the objective function is denoted
by $c_s: \mathcal{X}_s \rightarrow \mathbb{R}$ for a solution $x \in \mathcal{X}_s$.
Given a cost function $c_s(\cdot)$, the goal is to find the optimal solution $x^{s*}$ for
the given instance $s$:
\[
x^{s*} = \arg\min_{x \in \mathcal{X}_s} c_s(x).
\]

Because enumerating all possible solution sets for a CO problem is computationally
intractable, exact methods often relax the binary condition, solve a linear program (LP), and
subsequently branch on fractional solution components. This approach is known as branch and
bound (B\&B). While B\&B avoids enumerating entirely unnecessary solution sets, the number of
branches can still grow exponentially in the worst case.

\subsection{Traveling Salesman Problem}
\label{subsec:tsp}

Let $G = (V, E)$ be a complete undirected graph with $|V| = n$ cities and edge costs
$c_{ij} \ge 0$ for all $\{i,j\} \in E$. Define binary decision variables
\[
x_{ij} =
\begin{cases}
1, & \text{if the tour uses edge } \{i,j\},\\
0, & \text{otherwise.}
\end{cases}
\]

The traveling salesman problem (TSP) can be formulated as the following integer program:
\[
\begin{aligned}
\min_{x}\quad & \sum_{\{i,j\}\in E} c_{ij} x_{ij} \\
\text{s.t.}\quad
& \sum_{j\in V\setminus\{i\}} x_{ij} = 2, && \forall i\in V
  \quad \text{(degree constraints)}\\
& \sum_{\{i,j\}\in E(S)} x_{ij} \le |S|-1, && \forall S\subset V,\; 2\le |S|\le n-1
  \quad \text{(subtour elimination)}\\
& x_{ij}\in\{0,1\}, && \forall \{i,j\}\in E.
\end{aligned}
\]
Here, $E(S) = \{\{i,j\}\in E : i\in S,\; j\in S\}$ denotes the set of edges with both
endpoints in subset $S$. The subtour elimination constraints enforce that the selected edges
form a single Hamiltonian cycle visiting every city exactly once, rather than multiple
disjoint cycles.

\subsection{Capacitated Vehicle Routing Problem}
\label{subsec:cvrp}

Let $G=(V,E)$ be a complete graph where $V = \{0,1,\dots,n\}$, with depot $0$ and customers
$1,\dots,n$. Each customer $i\in\{1,\dots,n\}$ has demand $q_i>0$, and $c_{ij}\ge 0$ is the
travel cost between nodes $i$ and $j$. A fleet of at most $m$ identical vehicles is available,
each with capacity $Q$. Define binary decision variables
\[
x_{ij}=
\begin{cases}
1, & \text{if a vehicle travels directly from } i \text{ to } j,\\
0, & \text{otherwise,}
\end{cases}
\qquad \forall i\ne j,\; i,j\in V,
\]
and continuous load variables $u_i$ for $i\in\{1,\dots,n\}$ denoting the cumulative load
delivered upon arrival at customer $i$.

The capacitated vehicle routing problem (CVRP) can be formulated as the following
mixed-integer program:
\[
\begin{aligned}
\min_{x,u}\quad & \sum_{i\in V}\sum_{j\in V\setminus\{i\}} c_{ij}x_{ij}\\
\text{s.t.}\quad
& \sum_{j\in V\setminus\{i\}} x_{ij} = 1, && \forall i\in V\setminus\{0\}
  \quad \text{(each customer departs once)}\\
& \sum_{i\in V\setminus\{j\}} x_{ij} = 1, && \forall j\in V\setminus\{0\}
  \quad \text{(each customer arrives once)}\\
& \sum_{j\in V\setminus\{0\}} x_{0j} \le m, && \text{(number of routes)}\\
& \sum_{i\in V\setminus\{0\}} x_{i0}
  = \sum_{j\in V\setminus\{0\}} x_{0j}, && \text{(flow conservation at depot)}\\
& q_i \le u_i \le Q, && \forall i\in V\setminus\{0\}
  \quad \text{(capacity bounds)}\\
& u_i - u_j + Qx_{ij} \le Q - q_j, &&
  \forall i,j\in V\setminus\{0\},\; i\ne j
  \quad \text{(MTZ)}\\
& x_{ij}\in\{0,1\}, && \forall i\ne j,\; i,j\in V.
\end{aligned}
\]
The Miller--Tucker--Zemlin (MTZ) constraints enforce route connectivity and eliminate
subtours while simultaneously ensuring that vehicle capacity is not exceeded.

\subsection{Multi-Depot Vehicle Routing Problem}
\label{subsec:mdvrp}

Let $G=(V,E)$ be a complete graph containing multiple depots and customers. Let
$D=\{0,1,\dots,d\}$ be the set of depots and $C=\{d+1,\dots,d+n\}$ the set of customers, such
that $V=D\cup C$ and $D\cap C=\emptyset$. Depot $k\in D$ has at most $m_k$ identical vehicles
available, and each customer $i\in C$ has demand $q_i>0$. Let $c_{ij}\ge 0$ denote the travel
cost between nodes $i$ and $j$, and let $V_k=C\cup\{k\}$ be the customer set together with
depot $k$. Define depot-indexed binary decision variables
\[
x_{ij}^k=
\begin{cases}
1, & \text{if a vehicle from depot } k \text{ travels directly from } i \text{ to } j,\\
0, & \text{otherwise,}
\end{cases}
\qquad \forall k\in D,\; i\ne j,\; i,j\in V_k,
\]
and assignment variables
\[
y_{ik}=
\begin{cases}
1, & \text{if customer } i \text{ is served by depot } k,\\
0, & \text{otherwise,}
\end{cases}
\qquad \forall i\in C,\; k\in D.
\]

The multi-depot vehicle routing problem (MDVRP) can be formulated as the following
mixed-integer program:
\[
\begin{aligned}
\min_{x,y}\quad
& \sum_{k\in D}\sum_{i\in V_k}\sum_{j\in V_k\setminus\{i\}} c_{ij}x_{ij}^k\\
\text{s.t.}\quad
& \sum_{k\in D} y_{ik} = 1, && \forall i\in C
  \quad \text{(assign each customer to one depot)}\\
& \sum_{j\in V_k\setminus\{i\}} x_{ij}^k = y_{ik}, &&
  \forall i\in C,\; \forall k\in D
  \quad \text{(customer departure)}\\
& \sum_{j\in V_k\setminus\{i\}} x_{ji}^k = y_{ik}, &&
  \forall i\in C,\; \forall k\in D
  \quad \text{(customer arrival)}\\
& \sum_{j\in C} x_{kj}^k = \sum_{j\in C} x_{jk}^k \le m_k, && \forall k\in D
  \quad \text{(depot flow and route count)}\\
& \sum_{i\in S}\sum_{j\in S\setminus\{i\}} x_{ij}^k \le |S|-1, &&
  \forall k\in D,\; \forall S\subseteq C,\; 2\le |S|\le n-1
  \quad \text{(subtour elimination)}\\
& x_{ij}^k\in\{0,1\}, && \forall k\in D,\; i\ne j,\; i,j\in V_k,\\
& y_{ik}\in\{0,1\}, && \forall i\in C,\; k\in D.
\end{aligned}
\]
If vehicle capacities are included, standard capacity constraints, such as MTZ- or flow-based
constraints within each depot's routes, can be added to ensure that the total demand on any
route does not exceed the vehicle capacity.

\section{Classical and Learning-Based Solvers}
\label{sec:solvers}

Combinatorial optimization solvers can be broadly grouped into exact methods,
relaxation-based methods, and heuristic or metaheuristic methods. In practice, modern
state-of-the-art solvers frequently combine these ideas.

\subsection{Exact Solvers}
\label{subsec:exact-solvers}

Exact solvers provide guaranteed optimality, but their computational cost can grow
exponentially. Branch and bound (B\&B) builds a search tree over partial solutions. At each
node, it computes a lower bound, typically through a relaxation such as an LP relaxation. If
the bound is worse than the best-known feasible solution, which acts as an upper bound, the
node is pruned. B\&B is exact, but the search tree can still become exponentially large.

Branch and cut (B\&C) augments B\&B with cutting planes. When the LP relaxation yields a
fractional solution, the solver adds valid inequalities, or cuts, that remove the fractional
solution while preserving all integer-feasible solutions. For the TSP, subtour elimination
constraints serve as classic cuts. This strategy forms the backbone of solvers such as
Concorde \citep{applegate2006traveling}.

Branch and price (B\&P) combines B\&B with column generation. Instead of enumerating all
variables, which is often exponential, the solver optimizes a restricted master problem and
repeatedly generates new variables by solving a pricing subproblem. VRP formulations that use
route variables commonly employ this approach.

\subsection{Relaxation and Decomposition Methods}
\label{subsec:relaxation-decomposition}

Relaxation and decomposition methods are widely used to obtain bounds and exploit scalable
problem structure. LP relaxation replaces integrality constraints with continuous constraints
such as $x \in [0,1]$, yielding a linear program. LP relaxations provide lower bounds and are
standard components in B\&B and B\&C frameworks, although their fractional solutions may differ
substantially from integer-feasible solutions.

Lagrangian relaxation moves hard constraints into the objective function using multipliers.
This transformation can produce decomposable subproblems and often yields strong bounds,
especially for routing and scheduling problems. Benders decomposition instead splits variables
into a master problem and a subproblem. The subproblem generates Benders cuts that iteratively
refine the master problem, making the approach effective for problems with naturally separable
structure.

\subsection{Heuristics and Metaheuristics}
\label{subsec:heuristics-metaheuristics}

Heuristics and metaheuristics are designed to obtain high-quality feasible solutions quickly.
Local search iteratively improves a feasible solution through small neighborhood moves. In
routing problems, common operations include 2-opt, 3-opt, relocate, swap, and cross-exchange.
Although local search is fast and effective, it is susceptible to local minima.

Metaheuristics are higher-level strategies designed to escape local minima or explore the
solution space more broadly. Prominent examples include simulated annealing (SA), which
occasionally accepts worse moves according to a temperature schedule; tabu search, which uses
memory to avoid cycling back to recently visited solutions; genetic algorithms (GA), which
maintain a population of solutions through crossover and mutation; and ant colony optimization
(ACO), which constructs solutions probabilistically based on pheromone trails. These methods
do not guarantee optimality but are highly competitive for large instances.

\subsection{Learning-Assisted Solvers}
\label{subsec:learning-assisted-solvers}

Recent neural CO research often integrates traditional solvers in two distinct roles. First,
exact or high-quality heuristic solutions can serve as teachers or oracles, providing
ground-truth labels for supervised learning. Second, classical methods can act as decoding or
refinement steps: neural models generate a soft solution, such as a probabilistic heatmap,
which is subsequently transformed into a feasible hard solution using greedy decoding, local
search such as 2-opt, or other solver components.

In summary, exact solvers provide optimality guarantees at high computational cost,
relaxations and decompositions offer scalable bounds, and heuristics provide strong solutions
rapidly. Modern NCO systems frequently combine all three, aiming to address the limitations of
traditional solvers through data-driven advantages.

\section{Neural Network Methodologies}
\label{sec:neural-network-methodologies}

\subsection{Graph Neural Networks}
\label{subsec:gnn}

Graph neural networks (GNNs) are neural architectures designed to process data defined on
graphs by repeatedly performing message passing. Each node updates its representation by
aggregating transformed features from its neighbors. In each layer, a generic GNN computes
\[
h_i^{\ell+1}=\phi\Big(h_i^{\ell},\;
\operatorname{AGG}_{j\in\mathcal{N}(i)}\psi(h_i^{\ell},h_j^{\ell},e_{ij})\Big),
\]
where $h_i^{\ell}$ is the node embedding at layer $\ell$, $e_{ij}$ is an optional edge
feature, and $\operatorname{AGG}$ is a permutation-invariant operator such as summation or
mean aggregation.

In neural combinatorial optimization, classical backbones such as graph convolutional
networks (GCNs) \citep{kipf2016semi} and graph attention networks (GATs)
\citep{velickovic2017graph} primarily produce node embeddings. This is sufficient when the
target variable is node-based, such as predicting whether a node is selected in the maximum
independent set (MIS). However, for routing problems such as the TSP, the key latent variable
is often edge-based, indicating whether an edge belongs to the tour. Consequently, a standard
node-centric GNN requires additional structural mechanisms to reliably represent and predict
edge decisions.

\subsection{Anisotropic Graph Neural Networks}
\label{subsec:agnn}

Anisotropic graph neural networks (AGNNs) extend traditional message passing by maintaining
both node and edge embeddings and by using edge gates so that information flow depends on the
current edge state. This architecture is especially useful for diffusion-style denoising
networks, where the model ingests noisy variables $x_t$, such as node states for MIS or edge
states for TSP, and predicts the clean variables $x_0$.

Let $h_i^{\ell}$ and $e_{ij}^{\ell}$ denote node and edge features at layer $\ell$, and let
$t$ be a sinusoidal timestep embedding. A standard AGNN layer used in diffusion backbones
first performs an anisotropic edge update:
\[
\hat{e}_{ij}^{\ell+1}=P^{\ell}e_{ij}^{\ell}+Q^{\ell}h_i^{\ell}+R^{\ell}h_j^{\ell}.
\]
The edges are then refined with residual multilayer perceptrons (MLPs) and timestep
conditioning:
\[
e_{ij}^{\ell+1}
=e_{ij}^{\ell}+\operatorname{MLP}_e(\operatorname{BN}(\hat{e}_{ij}^{\ell+1}))
+\operatorname{MLP}_t(t).
\]

The node update then uses edge-gated messages:
\[
h_i^{\ell+1}=h_i^{\ell}+\alpha\Big(
\operatorname{BN}\big(U^{\ell}h_i^{\ell}
+\sum_{j\in\mathcal{N}(i)}
\sigma(\hat{e}_{ij}^{\ell+1})\odot V^{\ell}h_j^{\ell}\big)\Big).
\]
Here $\sigma(\cdot)$ is a sigmoid gate and $\odot$ is the Hadamard product. Intuitively, each
neighbor's message is filtered by the current edge embedding, allowing the network to
emphasize or suppress interactions in a direction- and edge-dependent manner.

Because AGNNs jointly model nodes and edges, they naturally support optimization tasks that
require edge-variable prediction, such as the TSP. For these problems, $e_{ij}^{0}$ is
initialized from the noisy edge variables $x_t$, and $h_i^{0}$ is set to positional or
sinusoidal node features. The final network head then predicts denoised edge probabilities or
values.

\section{Diffusion Models for COP}
\label{sec:diffusion-cop}

\subsection{DIFUSCO}
\label{subsec:difusco}

DIFUSCO (Diffusion for Combinatorial Optimization) solves CO problems by casting a feasible
solution as a binary structure, learning a diffusion-style denoising process to transform
random noise into that structure, and finally applying a classical decoder, such as greedy
decoding followed by 2-opt, to yield a valid tour or route \citep{sun2023difusco}.

\subsubsection{Representation}
\label{subsec:difusco-representation}

DIFUSCO represents the target solution as a binary variable $x_0\in\{0,1\}^N$. For TSP,
$x_0$ typically corresponds to edge decisions, indicating whether an edge is included in the
tour. Thus, $N$ scales with the number of candidate edges.

\subsubsection{Forward Process}
\label{subsec:difusco-forward}

A Markov chain gradually corrupts the optimal solution $x_0$ into random noise $x_T$ using a
variance schedule $\{\beta_t\}_{t=1}^T$. In discrete diffusion, the solution $x$ is converted
to a one-hot representation $\tilde{x}\in\{0,1\}^{N\times 2}$ and subjected to a categorical
transition:
\[
q(x_t\mid x_{t-1}) = \mathrm{Cat}(x_t;\,p=\tilde{x}_{t-1}Q_t),\quad
Q_t=\begin{pmatrix}
1-\beta_t & \beta_t\\
\beta_t & 1-\beta_t
\end{pmatrix}.
\]
As $\prod_{t=1}^T(1-\beta_t)\approx 0$, the terminal distribution becomes near-uniform.

For continuous diffusion on discrete data, the discrete solution $x_0\in\{0,1\}^N$ is
rescaled to $\hat{x}_0\in\{-1,1\}^N$ and subjected to Gaussian corruption:
\[
q(\hat{x}_t\mid \hat{x}_{t-1})
:= \mathcal{N}(\hat{x}_t;\sqrt{1-\beta_t}\,\hat{x}_{t-1},\,\beta_t I).
\]
This ensures that $\hat{x}_T$ approaches a standard Gaussian distribution when
$\prod_{t=1}^T(1-\beta_t)\approx 0$.

\subsubsection{Reverse Process}
\label{subsec:difusco-reverse}

DIFUSCO learns the reverse transitions defined by
\[
p_\theta(x_0) := \int p_\theta(x_{0:T})\,dx_{1:T},\quad
p_\theta(x_{0:T}) = p(x_T)\prod_{t=1}^T p_\theta(x_{t-1}\mid x_t).
\]
This is achieved by optimizing the standard variational lower bound, which yields
Kullback--Leibler (KL) divergence terms of the form
\[
D_{\mathrm{KL}}\big(q(x_{t-1}\mid x_t,x_0)\,\|\,p_\theta(x_{t-1}\mid x_t)\big).
\]

The neural network prediction target depends on the diffusion type. In discrete diffusion,
the denoiser is trained to predict the clean variable $p_\theta(x_0\mid x_t)$. The reverse
step is formed by substituting the prediction into the true posterior:
\[
p_\theta(x_{t-1}\mid x_t) \approx q(x_{t-1}\mid x_t,\hat{x}_0),\quad
\hat{x}_0\sim p_\theta(x_0\mid x_t).
\]
In continuous diffusion, the denoiser predicts the Gaussian noise $\varepsilon_t$:
\[
\varepsilon_t
= \frac{\hat{x}_t-\sqrt{\bar{\alpha}_t}\,\hat{x}_0}
{\sqrt{1-\bar{\alpha}_t}}
= f_\theta(\hat{x}_t,t),
\]
which is used to formulate a point estimate of $\hat{x}_0$ inside the Gaussian posterior for
$\hat{x}_{t-1}$.

\subsubsection{AGNN Denoiser}
\label{subsec:difusco-agnn-denoiser}

The conditional model $f_\theta(\cdot,t)$ is implemented using an edge-aware AGNN backbone
with timestep conditioning. After sampling the denoised representation to obtain a continuous
or probabilistic $x_0$, the model applies thresholding or quantization to map the heatmap back
to the $\{0,1\}$ domain. Finally, a non-differentiable decoder, such as greedy edge selection,
and optional local improvement heuristics, such as 2-opt, are applied to enforce hard
feasibility constraints and minimize the final objective value.

\section{Reinforcement Learning Fine-Tuning}
\label{sec:rl-ft}

\subsection{CADO}
\label{subsec:cado}

While diffusion-based models such as DIFUSCO demonstrate strong performance using supervised
learning (SL), this paradigm suffers from a fundamental objective mismatch. In SL, the neural
network is trained to imitate optimal ground-truth solutions by minimizing surrogate losses,
such as cross-entropy. However, minimizing this imitation loss does not guarantee minimizing
the actual decoded solution cost.

CADO (Cost-Aware Diffusion Solvers for Combinatorial Optimization through RL Fine-Tuning)
analyzes this objective mismatch through two main deficiencies \citep{yoon2024cado}. The
first is decoder-blindness: the SL objective is oblivious to the non-differentiable decoding
heuristic $f_g$, such as greedy search or 2-opt, that projects the continuous heatmap into a
discrete feasible solution. The second is cost-blindness: SL prioritizes structural similarity
to the ground-truth solution over the actual objective value of the problem. Empirical results
in CADO show that structural proximity to an optimal solution, measured by Hamming distance,
can have near-zero correlation with the final decoded solution cost.

To address these issues, CADO introduces a reinforcement learning (RL) fine-tuning framework
that explicitly aligns the diffusion process with the true CO objective.

\subsection{MDP Formulation}
\label{subsec:cado-mdp}

Instead of relying only on structural imitation, CADO directly optimizes the post-decoded
solution cost by formulating the iterative diffusion denoising process as a Markov decision
process (MDP). This formulation enables the model to be trained using the REINFORCE
algorithm.

Let the MDP be represented by the tuple $(\mathcal{S}, \mathcal{A}, P, \rho_0, R)$. The
state $s_t\in\mathcal{S}$, action $a_t\in\mathcal{A}$, and reward function $R(s_t,a_t)$ at
timestep $t$ are defined as
\[
s_t \triangleq (g, T-t, x_{T-t}), \quad a_t \triangleq x_{t-1},
\]
and
\[
R(s_t, a_t) \triangleq
\begin{cases}
-c_g(f_g(x_0)), & \text{if } t = T, \\
0, & \text{otherwise.}
\end{cases}
\]
Here, $g$ represents the problem instance, $c_g(\cdot)$ is the true cost function, and
$f_g(\cdot)$ is the non-differentiable decoder. By evaluating the reward only at the final
step, after the heatmap has been fully decoded, this formulation explicitly incorporates both
the decoder's behavior and the true cost feedback, resolving decoder-blindness and
cost-blindness.

\subsection{Reward Strategies}
\label{subsec:cado-reward-strategies}

To stabilize RL training and reduce variance, CADO uses two reward formulation strategies.
The standard reward (SR) uses the negative true solution cost directly as the reward signal.
This label-free formulation enables the model to learn through trial-and-error exploration and
can remain applicable even when labeled solutions are unavailable.

The label-centered reward (LCR) uses information from the SL pre-training dataset without
returning to pure imitation. It repurposes the ground-truth solution cost $b_D(g)$ as an
instance-specific baseline and defines the reward as the negative optimality gap:
\[
R_t = -\big(c_g(f_g(x_0)) - b_D(g)\big).
\]
Because the baseline $b_D(g)$ depends only on the problem instance and is independent of the
policy's actions, it serves as an unbiased baseline. Thus, the policy gradient remains
consistent with true cost minimization even if the dataset labels are computationally
suboptimal.

\subsection{Hybrid Fine-Tuning}
\label{subsec:cado-hybrid-ft}

Applying RL fine-tuning naively to the large AGNN architectures used in diffusion models can
be unstable and memory-intensive. To achieve parameter-efficient and stable adaptation, CADO
employs a hybrid fine-tuning (Hybrid-FT) strategy.

Hybrid-FT applies low-rank adaptation (LoRA) to the input layer and most backbone GNN layers.
This preserves the robust features learned during SL pre-training while introducing only a
small set of trainable low-rank matrices. At the same time, CADO applies selective layer
fine-tuning (Selective-FT), in which all parameters are retrained, to the final GNN layer and
the output layer. These layers are most directly responsible for final heatmap generation. The
combination supports rapid initial convergence while maintaining sufficient expressive power
for cost-aware adaptation.

